\markdownRendererDocumentBegin
\markdownRendererSectionBegin
\markdownRendererHeadingOne{Response to Reviewer \markdownRendererHash{}4}\markdownRendererInterblockSeparator
{}Manuscript: \markdownRendererEmphasis{Embedded Artificial Intelligence in Battery Management Systems: Pruning and Quantization for Efficient State-of-Charge and State-of-Health Estimation}\markdownRendererParagraphSeparator
{}The page and line references below refer to the final 62-page line-numbered anonymous revision compiled on 22 July 2026. Figure captions and other float contents are identified by figure number and page because the \markdownRendererCodeSpan{lineno} package does not assign reliable line numbers inside floats. No new battery experiment, model-training campaign, or STM32 measurement was performed for these comments. The two papers explicitly suggested by the reviewer were imported from the author's Zotero bibliography and are cited as Refs. [16] and [17]. Figures 2 and 7 retain the established visual language of the original manuscript with clarified captions, while Figure 6 provides the revised pruning schematic.\markdownRendererInterblockSeparator
{}\markdownRendererSectionBegin
\markdownRendererHeadingTwo{General response}\markdownRendererInterblockSeparator
{}We thank the reviewer for the detailed comments on presentation, literature context, and mathematical clarity. We shortened and refocused the Abstract, added and analysed the two suggested references, reorganised the method description around the actual Base, Pruned, and Quantized transformations, and added implementation-level equations for pruning, quantization, mixed-precision inference, and evaluation metrics. We also divided long paragraphs, performed a manuscript-wide language and logic pass, clarified the captions of the general design-of-experiments and quantization schematics, and revised the pruning schematic.\markdownRendererInterblockSeparator
{}
\markdownRendererSectionEnd \markdownRendererSectionBegin
\markdownRendererHeadingTwo{Comment 1}\markdownRendererInterblockSeparator
{}\markdownRendererBlockQuoteBegin
The expression of the abstract should be improved. The length should be suitable for the abstract, not too long.\markdownRendererBlockQuoteEnd \markdownRendererInterblockSeparator
{}\markdownRendererStrongEmphasis{Response:} We agree. The Abstract was rewritten close to the original narrative rather than as a list of qualifications. It introduces the engineering application and implemented artificial intelligence, establishes the full-precision estimators as the accuracy reference, and then presents the long-horizon and hardware-in-the-loop results. The reported flash and inference-time savings are tied explicitly to the evaluated 30\markdownRendererPercentSign{} pruning configuration, while the accompanying analytical model relates the computational reduction to pruning fraction and network dimensions. Detailed qualification of the timing-derived energy estimate and transferability is retained in the Methods and Discussion, where it can be explained without interrupting the Abstract. All acronyms are defined at first use. The title contains no acronym, and the keyword list has been reduced to six written-out terms.\markdownRendererParagraphSeparator
{}\markdownRendererStrongEmphasis{Changes in the manuscript:} revised Abstract and keywords (page 1, lines 6--26; Abstract lines 6--23 and keywords lines 24--26).\markdownRendererInterblockSeparator
{}
\markdownRendererSectionEnd \markdownRendererSectionBegin
\markdownRendererHeadingTwo{Comment 2}\markdownRendererInterblockSeparator
{}\markdownRendererBlockQuoteBegin
Please conduct more reference analysis for the research topic, such as An improved parameter identification and radial basis correction-differential support vector machine strategies for state-of-charge estimation of urban-transportation-electric-vehicle lithium-ion batteries, Improved singular filtering-Gaussian process regression-long short-term memory model for whole-life-cycle remaining capacity estimation of lithium-ion batteries adaptive to fast aging and multi-current variations, and so on.\markdownRendererBlockQuoteEnd \markdownRendererInterblockSeparator
{}\markdownRendererStrongEmphasis{Response:} We added both suggested articles from the author's Zotero export and analysed their relationship to the present work. The revised Related Work explains that the first combines online parameter identification, adaptive radial-basis correction, and a differential support-vector machine for SOC estimation under dynamic vehicle tests, while the second combines singular filtering, Gaussian-process regression, and an LSTM for whole-life capacity estimation under fast ageing and varying currents. We then state the relevant distinction: these studies broaden estimation methods and operating conditions, whereas neither reports the paired post-training pruning/quantization comparison with linked-firmware memory and recurrent microcontroller timing investigated here. They are therefore treated as complementary evidence, not as direct benchmark substitutes.\markdownRendererParagraphSeparator
{}\markdownRendererStrongEmphasis{Changes in the manuscript:} new comparative literature analysis and citations [16,17] (page 6, lines 143--154); full bibliography entries with DOI metadata (pages 57--58, lines 1022--1035).\markdownRendererInterblockSeparator
{}
\markdownRendererSectionEnd \markdownRendererSectionBegin
\markdownRendererHeadingTwo{Comment 3}\markdownRendererInterblockSeparator
{}\markdownRendererBlockQuoteBegin
The expression logic can be improved for your proposed method so that the innovation can be clearly understood. More mathematical analysis should be conducted to clarify the methods.\markdownRendererBlockQuoteEnd \markdownRendererInterblockSeparator
{}\markdownRendererStrongEmphasis{Response:} We reorganised the methods to distinguish three operations explicitly: the common FP32 Base equations, structured pruning that changes the hidden dimension and dense parameter submatrices, and quantization that changes the stored representation of selected recurrent weights without changing topology. The revised text now gives the analytical MAC expression and derives the general relative reduction for a pruning fraction \markdownRendererCodeSpan{p}, including its large-model limit \markdownRendererCodeSpan{2p-p\markdownRendererCircumflex{}2}. It then gives the retained-index submatrix construction in Eq. (17), the exact symmetric per-row quantizer in Eqs. (18)--(20), the reconstruction-error bound in Eq. (21), and the mixed-precision STM32 gate equation in Eq. (22). It also states which quantities remain FP32 and why storage reduction does not imply fully integer execution. These additions clarify that the contribution is an auditable paired deployment benchmark rather than a new recurrent cell or compression algorithm.\markdownRendererParagraphSeparator
{}\markdownRendererStrongEmphasis{Changes in the manuscript:} common architecture and general analytical MAC scaling (pages 15--16, lines 327--362); structured pruning definition and dense submatrix construction (pages 17--20, lines 387--434; Fig. 6 on page 20); exact quantization and mixed-precision equations (pages 20--22, lines 441--480; Fig. 7 on page 22); explicit benchmark objective and contributions (pages 4--5, lines 78--126).\markdownRendererInterblockSeparator
{}
\markdownRendererSectionEnd \markdownRendererSectionBegin
\markdownRendererHeadingTwo{Comment 4}\markdownRendererInterblockSeparator
{}\markdownRendererBlockQuoteBegin
Grammar should be checked and improved for the entire content. Please try to make every sentence correct and easy to understand.\markdownRendererBlockQuoteEnd \markdownRendererInterblockSeparator
{}\markdownRendererStrongEmphasis{Response:} We performed a manuscript-wide language pass and simplified ambiguous or overly compressed sentences. Particular attention was given to the Abstract, Introduction, Related Work, preprocessing definitions, pruning and quantization methods, Results interpretation, Discussion, Limitations, Conclusion, and declarations. We also corrected an author-name encoding error and removed duplicated acknowledgement text.\markdownRendererParagraphSeparator
{}\markdownRendererStrongEmphasis{Representative changes in the manuscript:} Abstract (page 1, lines 6--23); Introduction and Related Work (pages 4--8, lines 78--216); Methods (pages 9--27, lines 220--589); Results and Discussion (pages 27--40, lines 590--876); Conclusion and Outlook (pages 41--42, lines 877--917).\markdownRendererInterblockSeparator
{}
\markdownRendererSectionEnd \markdownRendererSectionBegin
\markdownRendererHeadingTwo{Comment 5}\markdownRendererInterblockSeparator
{}\markdownRendererBlockQuoteBegin
Please pay more attention to the expression logic and ensure the logic is correct for every sentence.\markdownRendererBlockQuoteEnd \markdownRendererInterblockSeparator
{}\markdownRendererStrongEmphasis{Response:} We agree and revised the argument flow from motivation to evidence boundary. The Introduction now closes with one objective and three contributions; Related Work ends with the precise research gap; the methods separate data, targets, common architecture, compression transformations, and measurement definitions; the Results distinguish observations from causal claims; and the Discussion separates measured findings, implementation-based interpretations, and unevaluated transferability. Unsupported statements, including a causal regularisation explanation for the lower Pruned SOC MAE and independent energy-measurement wording, were removed or bounded.\markdownRendererParagraphSeparator
{}\markdownRendererStrongEmphasis{Changes in the manuscript:} objective and contribution sequence (page 5, lines 111--126); research gap and implementation boundary (pages 7--8, lines 168--216); Base/Pruned/Quantized method boundary (pages 15--22, lines 327--480); evidence-based interpretation and transferability discussion (pages 37--40, lines 761--876); Conclusion and Outlook (pages 41--42, lines 877--917).\markdownRendererInterblockSeparator
{}
\markdownRendererSectionEnd \markdownRendererSectionBegin
\markdownRendererHeadingTwo{Comment 6}\markdownRendererInterblockSeparator
{}\markdownRendererBlockQuoteBegin
Some paragraphs are too long. Please optimize your expression.\markdownRendererBlockQuoteEnd \markdownRendererInterblockSeparator
{}\markdownRendererStrongEmphasis{Response:} We split and tightened long paragraphs throughout the manuscript. The revised structure separates dataset design from check-up procedures, training settings from deployment state handling, benchmark replay from firmware measurement, static operation counts from timing interpretation, and demonstrated limitations from future work. This reduces sentence load while preserving the quantitative details required for reproducibility.\markdownRendererParagraphSeparator
{}\markdownRendererStrongEmphasis{Representative changes in the manuscript:} dataset and preprocessing (pages 9--15, lines 220--326); training and deployment descriptions (pages 15--17, lines 327--386); benchmark methodology (pages 23--27, lines 481--589); Discussion and transferability (pages 37--40, lines 761--876); Conclusion and Outlook (pages 41--42, lines 877--917).\markdownRendererInterblockSeparator
{}
\markdownRendererSectionEnd \markdownRendererSectionBegin
\markdownRendererHeadingTwo{Comment 7}\markdownRendererInterblockSeparator
{}\markdownRendererBlockQuoteBegin
Some figures are unclear. Please try to optimize them.\markdownRendererBlockQuoteEnd \markdownRendererInterblockSeparator
{}\markdownRendererStrongEmphasis{Response:} We revised the presentation of three explanatory figures. Figure 2 uses the general three-factor design-of-experiments geometry from the original manuscript, while its caption now explains the generic condition labels and directs the reader to the dataset-specific factor levels in the text. Figure 6 depicts the dense hidden-channel pruning path and the implemented SOC and SOH dimension changes. Figure 7 restores the more detailed FP32-to-INT8 deployment schematic from the original manuscript. Its caption clarifies that the INT8 data type spans \markdownRendererCodeSpan{-128...127}, while the implemented symmetric quantizer generates \markdownRendererCodeSpan{-127...127} and leaves \markdownRendererCodeSpan{-128} unused. It also states that only the recurrent matrices are quantized and that the remaining operations stay in FP32.\markdownRendererParagraphSeparator
{}\markdownRendererStrongEmphasis{Changes in the manuscript:} restored general Fig. 2 with clarified caption (page 10); revised Fig. 6 and caption (page 20); restored detailed Fig. 7 with implementation-consistent caption (page 22).\markdownRendererInterblockSeparator
{}
\markdownRendererSectionEnd \markdownRendererSectionBegin
\markdownRendererHeadingTwo{Comment 8}\markdownRendererInterblockSeparator
{}\markdownRendererBlockQuoteBegin
More expression about the mathematical analysis should be conducted to show your idea clearly. Please try to highlight your proposed method and focus on it.\markdownRendererBlockQuoteEnd \markdownRendererInterblockSeparator
{}\markdownRendererStrongEmphasis{Response:} In addition to the method equations described in our response to Comment 3, we formalised the evaluation criteria so the relation between predictions and reported percentages is unambiguous. The revised manuscript defines signed percentage-point error, MAE, RMSE, and the empirical P95 threshold in Eq. (24) and the following expression. The pruning analysis now derives the finite-model reduction as a function of both hidden size and pruning fraction and proves the large-model limit \markdownRendererCodeSpan{2p-p\markdownRendererCircumflex{}2}. Figure A.21 visualises this relation from 0\markdownRendererPercentSign{} to 60\markdownRendererPercentSign{} pruning and highlights the evaluated 30\markdownRendererPercentSign{} point, whose asymptotic limit is 51\markdownRendererPercentSign{}. It also distinguishes these architecture-level counts from measured flash, latency, and the timing-derived energy proxy. This mathematical treatment focuses the contribution on a transparent, reproducible comparison of two deployment transformations.\markdownRendererParagraphSeparator
{}\markdownRendererStrongEmphasis{Changes in the manuscript:} analytical MAC count and general pruning-fraction reduction (pages 15--16, lines 335--351); pruning equations and operating-point scope (pages 18--19, lines 394--434); quantization and reconstruction equations (pages 21--22, lines 448--472); formal error metrics (pages 25--26, lines 559--564); measured/static interpretation (page 38, lines 784--813); analytical scaling in Appendix A.6 (page 50) and Fig. A.21 (page 51).
\markdownRendererSectionEnd 
\markdownRendererSectionEnd \markdownRendererDocumentEnd